% !Mode:: "TeX:UTF-8"

%%% 说明: 此部分需要自己填写的内容:  (1) 中文摘要及关键词 (2) 英文摘要及关键词
%--------------------------------------------------------------------------
%%%页眉设置
%
%\pagestyle{fancyplain}
%\fancyhf{}  %清除以前对页眉页脚的设置
%\renewcommand{\headrulewidth}{0.5pt}
%\fancyfoot[C]{-\,\thepage\,-}
%\fancypagestyle{plain}
%{
%\fancyhead[CE]{武汉大学博士学位论文}
%\fancyhead[CO]{}
%\renewcommand{\headrulewidth}{0.5pt}
%\fancyfoot{}                                    % clear all footer fields
%\fancyfoot[C]{-\,\thepage\,-}
%}

%%%%%%%%%%%%%%%%%%%%%%%
%%% ------------ 中文摘要 ---------------%%%
%%%%%%%%%%%%%%%%%%%%%%%
\begin{cnabstract}

偏微分方程（PDEs）是物理学、工程学等领域核心的数学建模工具。近年来，以深度里兹法（DRM）和物理信息神经网络（PINNs）为代表的深度学习方法在科学计算领域取得了重要进展，为克服传统网格类数值方法在高维空间中遭遇的“维数灾难”提供了新的思路。然而，在处理特征值正交约束、强非线性分岔以及反问题不适定性等复杂物理场景时，现有神经网络求解器仍面临优化瓶颈与理论支撑不足的问题。针对上述关键问题，本文以复杂偏微分系统为研究对象，开展了一系列基于神经网络的建模与计算方法研究。本文的主要工作与研究成果如下：

第一，针对高维椭圆型多重特征值问题，提出了一种基于惩罚变分原理的无约束DRM求解框架。传统神经网络方法面临特征函数正交约束难以实现且缺乏泛化理论支撑的困境。本文通过引入软边界惩罚与正交化约束，实现了多阶特征值的同步求解。在理论层面，本文建立了一套误差分解体系，刻画了特征函数的内在正则性，并推导了由网络深度、宽度及样本量显式表达的总误差上界，在特征值充分光滑的假设下，算法最优可达 $\mathcal{O}(n^{-1/6})$ 的收敛率，从理论上表明了其缓解维数灾难的潜力。高达20维区域的数值实验，验证了该方法在高维场景下的可靠性。

其次，针对带有 Dirichlet 和 Robin 边界条件的半线性椭圆方程，构建了基于惩罚变分形式的深度学习求解模型。利用具有 $\mathrm{ReLU}^2$ 激活函数的残差神经网络，证明了变分问题极小值点的存在性与唯一性。本文推导了数值解与真实解之间 $H^1$ 范数误差的显式上界，该上界由逼近误差、统计误差和截断误差组成。理论分析表明，随着网络规模的增加和样本量的增大，数值解能够以确定的速率收敛于真实解。数值实验验证了该方法在高维及非光滑解情形下的较好精度表现。

最后，针对椭圆方程势函数重构的反问题，提出了结合 Tikhonov 正则化的PINNs方法。该方法通过在损失函数中引入 $H^2$ 范数正则化项，缓解了反问题的不适定性。本文建立了反势能问题的稳定性估计，并利用统计学习理论推导了超额风险界。理论结果表明了重构解关于噪声水平的收敛率。数值对比实验表明，与传统有限元方法相比，该方法在处理不连续势函数和噪声测量数据时具有较好的鲁棒性和较高的精度。

\end{cnabstract}
\vspace{1em}\par\vfill

%%%--------- 关键词 -------- %%%
\cnkeywords{偏微分方程，深度神经网络，特征值问题，反问题}

%%%%%%%%%%%%%%%%%%%%%%%


%%%%%%%%%%%%%%%%%%%%%%%
%%% ------------ 英文摘要 ---------------%%%
%%%%%%%%%%%%%%%%%%%%%%%

\begin{enabstract}
Partial differential equations (PDEs) serve as core mathematical modeling tools in fields such as physics and engineering. In recent years, deep learning methods—notably the Deep Ritz Method (DRM) and Physics-Informed Neural Networks (PINNs)—have achieved significant advancements in scientific computing. They provide a new approach for overcoming the "curse of dimensionality" that plagues traditional mesh-based numerical methods in high-dimensional spaces. However, when addressing complex physical scenarios involving the orthogonal constraints of eigenvalues, strong nonlinear bifurcations, and ill-posed inverse problems, existing neural network solvers still encounter optimization bottlenecks and lack sufficient theoretical foundations. To address these key challenges, this paper focuses on complex PDE systems. Centered on the core trajectory of "algorithm design tailored for complex problems, quantitative theoretical analysis, and numerical empirical validation," this work investigates neural network-based modeling and computational methodologies. The primary contributions and findings of this research are summarized as follows:

First, targeting high-dimensional elliptic multiple eigenvalue problems, we propose an unconstrained Deep Ritz framework based on a penalized variational principle. Traditional neural network approaches face the difficulty of implementing orthogonal constraints for eigenfunctions and the absence of theoretical support for generalization. By incorporating soft boundary penalties and orthogonalization constraints, this study achieves the synchronous and stable computation of multi-order eigenvalues. Theoretically, we establish an error decomposition framework that characterizes the intrinsic regularity of eigenfunctions. Furthermore, we derive an explicit upper bound for the total error expressed in terms of network depth, width, and sample size. Under the assumption of sufficient smoothness of the eigenfunctions, the algorithm attains an optimal convergence rate of $\mathcal{O}(n^{-1/6})$, thereby theoretically demonstrating its potential for alleviating the curse of dimensionality. Numerical experiments conducted on complex domains up to 20 dimensions (e.g., unit spheres) validate the accuracy and reliability of the proposed method in high-dimensional and extreme scenarios.

Second, we construct a deep learning solver modeled on a penalized variational formulation for semilinear elliptic equations subject to Dirichlet and Robin boundary conditions. Utilizing Residual Neural Networks (ResNets) with a $\mathrm{ReLU}^2$ activation function, we prove the existence and uniqueness of the minimizer for the variational problem. We derive an explicit upper bound for the $H^1$-norm error between the numerical and exact solutions, which comprises approximation error, statistical error, and truncation error. Theoretical analyses demonstrate that as the network scale and sample size increase, the numerical solution converges to the exact solution at a deterministic rate. Numerical experiments confirm the good performance of this method, particularly in scenarios involving high-dimensional and non-smooth solutions.

Finally, addressing the inverse problem of potential function reconstruction in elliptic equations, we propose a Physics-Informed Neural Network method integrated with Tikhonov regularization. This approach mitigates the ill-posed nature of the inverse problem by introducing an $H^2$-norm regularization term into the loss function. We establish a stability estimate for the inverse potential problem and employ statistical learning theory to derive the excess risk bounds. Theoretical outcomes indicate the convergence rate of the reconstructed solution relative to the noise level. Comparative numerical experiments demonstrate that, relative to traditional finite element methods, the proposed approach exhibits good robustness and higher accuracy when handling discontinuous potential functions and noisy measurement data.
\end{enabstract}
\vspace{1em}\par\vfill

%%%------ 英文关键词 ------- %%%
\enkeywords{partial differential equations, deep neural networks, error analysis, convergence rate}


